\documentclass[11pt,a4paper]{article}

% ---------------------------------------------------------------------------
% Packages
% ---------------------------------------------------------------------------
\usepackage[utf8]{inputenc}
\usepackage[T1]{fontenc}
\usepackage{amsmath,amssymb}
\usepackage{physics}
\usepackage{braket}
\usepackage{hyperref}
\usepackage[margin=1in]{geometry}
\usepackage{graphicx}
\usepackage{enumitem}
\usepackage{booktabs}
\usepackage{multirow}
\usepackage{xcolor}
\usepackage{listings}

% ---------------------------------------------------------------------------
% Custom commands
% ---------------------------------------------------------------------------
\newcommand{\Gammamatrix}[1]{\Gamma_{#1}}
\newcommand{\gfive}{\gamma_5}
\newcommand{\gsrc}{\Gamma_{\text{src}}}
\newcommand{\gsink}{\Gamma_{\text{sink}}}
\newcommand{\gseq}{\Gamma_{\text{seq}}}
\newcommand{\gcur}{\Gamma}
\newcommand{\Sq}{S_{q}}
\newcommand{\Santi}{S_{\text{anti}}}
\newcommand{\Sseq}{S_{\text{seq}}}
\newcommand{\etaseq}{\eta_{\text{seq}}}
\newcommand{\einsum}{\texttt{einsum}}
\newcommand{\pyquda}{\texttt{PyQUDA}}
\newcommand{\quda}{\texttt{QUDA}}
\newcommand{\hdfcmd}{\texttt{HDF5}}

\title{\textbf{Pion Electromagnetic Form Factor}\\[4pt]
\large Physics Note --- \pyquda{} Implementation}
\author{}
\date{\today}

\begin{document}
\maketitle

% ===========================================================================
\section{Physics Overview}
% ===========================================================================

The electromagnetic form factor (EMFF) of the pion, \(F_\pi(Q^2)\), parametrizes
the hadronic matrix element of the electromagnetic current between pion states.
For a charged pion \(\pi^+\), the matrix element is

\begin{equation}
\braket{\pi^+(p_f) | J_\mu^{\text{em}}(0) | \pi^+(p_i)}
  = (p_i + p_f)_\mu \, F_\pi(Q^2) ,
\qquad Q^2 = -q^2 = -(p_f - p_i)^2 \ge 0 .
\label{eq:emff_def}
\end{equation}

The electromagnetic current for a single quark flavor \(f\) with electric charge
\(e_f\) is \(J_\mu^{\text{em}} = e_f \, \bar\psi \gamma_\mu \psi\).  On the
lattice this is accessed through the Euclidean three-point correlation function

\begin{equation}
C_3^\Gamma(\mathbf{q}, \tau; \mathbf{p}_f, t_{\text{sep}})
  = \sum_{\mathbf{x}} \sum_{\mathbf{y}}
    e^{-i\mathbf{q}\cdot(\mathbf{x}-\mathbf{x}_0)}
    e^{-i\mathbf{p}_f\cdot(\mathbf{y}-\mathbf{x}_0)} \,
    \big\langle
      O_\pi(\mathbf{y}, t_{\text{sep}}) \,
      J^\Gamma(\mathbf{x}, \tau) \,
      O_\pi^\dagger(\mathbf{x}_0, 0)
    \big\rangle ,
\label{eq:C3_operator}
\end{equation}

where \(O_\pi = \bar\psi \gsink \psi\) is the pion interpolating operator, and
\(J^\Gamma = \bar\psi \gcur \psi\) is the local current insertion.  The initial
pion momentum is fixed by momentum conservation:
\(\mathbf{p}_i = \mathbf{p}_f - \mathbf{q}\).

The pion interpolators project onto definite three-momentum via Fourier
transforms at the sink (momentum \(\mathbf{p}_f\)) and the current insertion
(momentum transfer \(\mathbf{q}\)).  For the EMFF channel we are interested in
the vector components of the current (\(\gcur = \gamma_\mu\) with \(\mu = X, Y,
Z, T\)), but the implementation described here scans all 16 Dirac bilinears for
diagnostic completeness and cross-checking.

After Wick contraction, the connected three-point function takes the trace form

\begin{equation}
C_3^\Gamma(\mathbf{q}, \tau; \mathbf{p}_f, t_{\text{sep}})
  = \sum_{\mathbf{x}} \sum_{\mathbf{y}}
    e^{-i\mathbf{q}\cdot(\mathbf{x}-\mathbf{x}_0)}
    e^{-i\mathbf{p}_f\cdot(\mathbf{y}-\mathbf{x}_0)} \,
    \tr\Big[
      \Santi(\mathbf{y}, t_{\text{sep}}; \mathbf{x}_0, 0) \,
      \gsink \;
      \Sq(\mathbf{y}, t_{\text{sep}}; \mathbf{x}, \tau) \,
      \gcur \;
      \Sq(\mathbf{x}, \tau; \mathbf{x}_0, 0) \,
      \gsrc
    \Big] .
\label{eq:C3_trace}
\end{equation}

Here \(\Sq(x, y)\) denotes the quark propagator from spacetime point \(y\) to
\(x\), and \(\Santi(x, y) = \gfive \, \Sq(x, y)^\dagger \, \gfive\) is the
antiquark propagator obtained via \(\gfive\)-hermiticity.  The trace runs over
spin and color indices.  Equation~\eqref{eq:C3_trace} is the connected
contribution; disconnected diagrams (where the current couples to a sea-quark
loop) are not treated in this module.


% ===========================================================================
\section{Correlation Functions in Detail}
% ===========================================================================

\subsection{Propagator Layout Convention}

The \pyquda{} \texttt{LatticePropagator} stores data in a nine-dimensional
array with index layout

\begin{equation}
\texttt{prop.data} \sim (w, t, z, y, x, \alpha, \beta, c, d) ,
\label{eq:data_layout}
\end{equation}

where
\begin{itemize}[nosep]
  \item \(w \in \{0,1\}\) is the checkerboard parity (even/odd site);
  \item \((t, z, y, x)\) are spacetime coordinates (note the fine coordinate
        \(x\) is the last spatial index, reflecting the checkerboard
        decomposition);
  \item \(\alpha\) is the spin index at the sink (final point);
  \item \(\beta\) is the spin index at the source (initial point);
  \item \(c\) is the color index at the sink;
  \item \(d\) is the color index at the source.
\end{itemize}

In prose indices throughout this note (and in the \einsum{} patterns) we label
these as \((w, t, z, y, x, \alpha, \beta, a, b)\) with \(\alpha\) = spin\_sink,
\(\beta\) = spin\_src, \(a\) = color\_sink, \(b\) = color\_src, or with
equivalent index letters depending on context.

\subsection{Three-Point Function}

The connected three-point function, written with explicit spin and color
indices, is

\begin{equation}
\begin{aligned}
C_3^\Gamma(\mathbf{q}, \tau; \mathbf{p}_f, t_{\text{sep}})
  = \sum_{\mathbf{x},\mathbf{y}} &
    e^{-i\mathbf{q}\cdot(\mathbf{x}-\mathbf{x}_0)}
    e^{-i\mathbf{p}_f\cdot(\mathbf{y}-\mathbf{x}_0)} \\
  &\times \Santi\big|_{\mathbf{y},t_{\text{sep}}}^{\alpha\beta;ab}
    \; (\gsink)^{\beta\gamma} \;
    \Sq\big|_{\mathbf{y},t_{\text{sep}};\,\mathbf{x},\tau}^{\gamma\delta;bc}
    \; (\gcur)^{\delta\epsilon} \;
    \Sq\big|_{\mathbf{x},\tau;\,\mathbf{x}_0,0}^{\epsilon\zeta;ca}
    \; (\gsrc)^{\zeta\alpha} ,
\end{aligned}
\label{eq:C3_index}
\end{equation}

where repeated indices are summed over.  The color trace closes because
\(\Santi\) and \(\Sq\) have opposite color flow directions, yielding
\({\Santi}_{ab} \, {\Sq}_{bc} \, {\Sq}_{ca}\) --- a proper color singlet.

The intermediate antiquark-to-quark connection at the current insertion point
\((\mathbf{x}, \tau)\) is given by the standard quark propagator \(\Sq(\mathbf{y},
t_{\text{sep}}; \mathbf{x}, \tau)\), representing the propagation from the
current insertion up to the sink.

\subsection{Two-Point Function}

The pion two-point function, serving as a normalization and diagnostic, is

\begin{equation}
C_2^\Gamma(\mathbf{p}, t)
  = \sum_{\mathbf{x}}
    e^{-i\mathbf{p}\cdot(\mathbf{x}-\mathbf{x}_0)} \,
    \tr\Big[
      \Santi(\mathbf{x}, t; \mathbf{x}_0, 0) \,
      \gsink \;
      \Sq(\mathbf{x}, t; \mathbf{x}_0, 0) \,
      \gsrc
    \Big] .
\label{eq:C2_trace}
\end{equation}

For the standard pseudoscalar pion with \(\gsrc = \gsink = \gfive\), this
reduces to

\begin{equation}
C_2^{\gfive}(\mathbf{p}, t)
  = \sum_{\mathbf{x}}
    e^{-i\mathbf{p}\cdot(\mathbf{x}-\mathbf{x}_0)} \,
    \tr\Big[
      \gfive \, \Sq(\mathbf{x}, \mathbf{x}_0)^\dagger \, \gfive \;
      \gfive \;
      \Sq(\mathbf{x}, \mathbf{x}_0) \;
      \gfive
    \Big]
  = \sum_{\mathbf{x}}
    e^{-i\mathbf{p}\cdot(\mathbf{x}-\mathbf{x}_0)} \,
    \tr\Big[ \Sq(\mathbf{x}, \mathbf{x}_0)^\dagger \, \Sq(\mathbf{x}, \mathbf{x}_0) \Big] .
\label{eq:C2_g5}
\end{equation}


% ===========================================================================
\section{Sequential Source Method}
% ===========================================================================

The double sum over sink position \(\mathbf{y}\) in
Eq.~\eqref{eq:C3_index} is computationally prohibitive if performed by nested
loops.  The standard solution is the \emph{fixed-sink sequential source method},
which absorbs the sink sum into a single sequential inversion.

\subsection{Construction of the Sequential Source}

The sequential source is built on time slice \(t_{\text{sep}} = t_{\text{src}} +
t_{\text{insert}}\) (modulo \(L_t\)).  Starting from the (anti)quark propagator
\(\Sq^{\text{neg}}(\mathbf{y}, t_{\text{sep}}; \mathbf{x}_0, 0)\) with
appropriate sink smearing applied, the source is

\begin{equation}
\etaseq(\mathbf{y}; \mathbf{p}_f, t_{\text{sep}})
  = \delta_{t_y,\,t_{\text{sep}}} \;
    e^{-i\mathbf{p}_f\cdot(\mathbf{y}-\mathbf{x}_0)} \;
    \gseq \;
    \Sq^{\text{neg}}(\mathbf{y}, t_{\text{sep}}; \mathbf{x}_0, 0) ,
\label{eq:etaseq}
\end{equation}

where the sequential gamma structure is

\begin{equation}
\gseq = \gfive \, \gsink^\dagger \, \gfive .
\label{eq:gseq}
\end{equation}

The time-slicing operation \(\delta_{t_y,t_{\text{sep}}}\) is handled by the
utility function \texttt{sequential12}, which extracts the propagator data
restricted to the specified time slice.

For smeared-smeared three-point functions the code applies the same sink-side
smearing pattern to the active sequential line before the sequential inversion.
If \(\mathcal{S}_{w,\mathbf{k}}\) denotes the boosted Gaussian smearing
operator defined in Sec.~\ref{sec:boosted_smearing}, the actual right-hand side
used by the current production workflow is
\begin{equation}
\etaseq^{SS}(\mathbf{y}; \mathbf{p}_f, t_{\text{sep}})
  =
  \mathcal{S}_{w,\mathbf{k}_{\rm pos,sink}}
  \left[
    \delta_{t_y,\,t_{\text{sep}}} \;
    e^{-i\mathbf{p}_f\cdot(\mathbf{y}-\mathbf{x}_0)} \;
    \gseq \;
    \Sq^{\text{neg},\,S\to S}
      (\mathbf{y}, t_{\text{sep}}; \mathbf{x}_0, 0)
  \right] .
\label{eq:etaseq_smeared}
\end{equation}
Here \(\Sq^{\text{neg},\,S\to S}\) already contains source smearing and the
spectator-line sink smearing with \(\mathbf{k}_{\rm neg,sink}\), while the
outer \(\mathcal{S}_{w,\mathbf{k}_{\rm pos,sink}}\) supplies the missing
sink-side smearing for the active quark line represented by the sequential
propagator.  This is the meson analogue of the proton fixed-sink sequential
source convention.

\subsection{Sequential Inversion}

The sequential propagator \(\Sseq\) is obtained by solving the Dirac equation
with \(\etaseq\) as the source:

\begin{equation}
D(x, y) \; \Sseq(y, x_0; \mathbf{p}_f, t_{\text{sep}}) = \etaseq(x; \mathbf{p}_f, t_{\text{sep}}) .
\label{eq:D_Sseq}
\end{equation}

This inversion is performed using the same Dirac operator employed for the
forward propagators.  In the code, the call is
\(\texttt{core.invertPropagator}(\texttt{dirac}, \texttt{src\_seq}, 1, 0)\).

\subsection{Antiquark Conversion}

After the sequential inversion, the resulting propagator is converted into an
antiquark backward line using \(\gfive\)-hermiticity, exactly as is done for the
two-point function:

\begin{equation}
\Sseq^{\text{anti}}(x, x_0; \mathbf{p}_f, t_{\text{sep}})
  = \gfive \, \Sseq(x, x_0; \mathbf{p}_f, t_{\text{sep}})^\dagger \, \gfive .
\label{eq:Sseq_anti}
\end{equation}

\subsection{Sink-Summed Three-Point Function}

With the sequential antiquark propagator in hand, the double sum over \(\mathbf{y}\)
has been absorbed into the sequential source construction.  The three-point
function reduces to a single spatial sum over the current insertion point
\(\mathbf{x}\):

\begin{equation}
\boxed{
C_3^\Gamma(\mathbf{q}, \tau; \mathbf{p}_f, t_{\text{sep}})
  = \sum_{\mathbf{x}}
    e^{-i\mathbf{q}\cdot(\mathbf{x}-\mathbf{x}_0)} \,
    \tr\Big[
      \Sseq^{\text{anti}}(\mathbf{x}, \mathbf{x}_0; \mathbf{p}_f, t_{\text{sep}}) \,
      \gcur \;
      \Sq(\mathbf{x}, \tau; \mathbf{x}_0, 0) \,
      \gsrc
    \Big] .
}
\label{eq:C3_seq}
\end{equation}

This is the expression evaluated by the \texttt{contract\_EMFF} method in the
code.


% ===========================================================================
\section{Two-Point Function (Implementation)}
% ===========================================================================

The two-point contraction implemented in \texttt{contract\_2pt\_pion} follows
Eq.~\eqref{eq:C2_trace}.  The procedure is:

\begin{enumerate}[nosep]
  \item Apply boosted smearing at the sink to both \(\Sq\) (positive) and
        \(\Sq^{\text{neg}}\) (negative) propagators.
  \item Convert the negative propagator to an antiquark backward line:
        \(\Santi = \gfive (\Sq^{\text{neg}})^\dagger \gfive\).
  \item Insert the sink gamma on the backward line.
  \item Contract the backward line, forward line, and source gamma into a
        spin-color trace.
  \item Project onto definite momentum \(\mathbf{p}\) using the Fourier phase.
\end{enumerate}

The contribution from the sink-summed three-point signal is stored alongside
the two-point correlators in the \hdfcmd{} output tree under
\(\texttt{data/c2pt/}\).


% ===========================================================================
\section{Source Gamma Conventions}
\label{sec:src_gamma}
% ===========================================================================

The source gamma matrix \(\gsrc\) in Eq.~\eqref{eq:C3_trace} can be specified
through four modes, controlled by the parameter \texttt{src\_gamma}:

\begin{enumerate}[label=\textbf{Mode \arabic*:}, leftmargin=*, nosep]
  \item \texttt{"fixed\_g5"} --- The source gamma is fixed to \(\gfive\)
        regardless of the sink gamma.  This is the standard choice for
        pseudoscalar pion measurements, where \(\gsrc = \gfive\) projects onto
        the pion's \(\gfive\) quantum numbers from the source.

  \item \texttt{"same\_as\_sink"} --- The source gamma is set equal to the sink
        gamma for each gamma channel separately:
        \(\gsrc = \gsink\).  This is useful for exploring the full Dirac
        structure of the correlation functions across all 16 bilinears.

  \item \texttt{"dagger\_of\_sink"} --- The source gamma is the
        \(\gfive\)-hermitian conjugate of the sink gamma:
        \(\gsrc = \gfive \, \gsink^\dagger \, \gfive\).  This mode is
        motivated by the sequential source construction, where the same
        transformation appears.

  \item \textbf{Any gamma label} --- A specific gamma matrix label (e.g.,
        \texttt{"5"}, \texttt{"T"}, \texttt{"X"}, etc.) can be provided directly,
        fixing \(\gsrc\) to that single gamma structure for all channels.
\end{enumerate}

The implementation is in the function \texttt{\_source\_gamma\_stack}, which
takes as input the precomputed stack of 16 sink (or current) gamma matrices and
returns the corresponding stack of source gamma matrices according to the chosen
mode.


% ===========================================================================
\section{Gamma Matrix Convention}
% ===========================================================================

The code scans all 16 independent Dirac bilinears.  The gamma matrices are
labeled by the following mnemonics, ordered in the standard project convention:

\begin{center}
\begin{tabular}{llll}
\toprule
\textbf{Index} & \textbf{Label} & \textbf{Dirac Structure} & \textbf{QUDA Bitmask} \\
\midrule
0  & 5    & \(\gfive\)                  & 15 \\
1  & T    & \(\gamma_4\) (temporal)     & 8  \\
2  & T5   & \(\gamma_4 \gfive\)        & 7  \\
3  & X    & \(\gamma_1\) (spatial)     & 1  \\
4  & X5   & \(\gamma_1 \gfive\)        & 14 \\
5  & Y    & \(\gamma_2\) (spatial)     & 2  \\
6  & Y5   & \(\gamma_2 \gfive\)        & 13 \\
7  & Z    & \(\gamma_3\) (spatial)     & 4  \\
8  & Z5   & \(\gamma_3 \gfive\)        & 11 \\
9  & I    & \(\mathbb{1}\) (identity)  & 0  \\
10 & SXT  & \(\sigma_{xt}\)            & 9  \\
11 & SXY  & \(\sigma_{xy}\)            & 3  \\
12 & SXZ  & \(\sigma_{xz}\)            & 5  \\
13 & SYT  & \(\sigma_{yt}\)            & 10 \\
14 & SYZ  & \(\sigma_{yz}\)            & 6  \\
15 & SZT  & \(\sigma_{zt}\)            & 12 \\
\bottomrule
\end{tabular}
\end{center}

The \quda{} bitmask convention uses a 4-bit encoding where bit 0 corresponds to
\(\gamma_x\), bit 1 to \(\gamma_y\), bit 2 to \(\gamma_z\), and bit 3 to
\(\gamma_t\).  Gamma products are formed by XOR of the constituent bit patterns.
For the identity, the bitmask is 0; \(\gfive = \gamma_x \gamma_y \gamma_z
\gamma_t\) corresponds to \(1+2+4+8 = 15\).  The tensor bilinears
\(\sigma_{\mu\nu} = \frac{1}{2}[\gamma_\mu, \gamma_\nu]\) correspond to the XOR
of the individual bitmasks.

In the code, the gamma matrices are instantiated via \(\texttt{gamma.gamma}(n)\)
where \(n\) is the \quda{} bitmask index, and stored in a stacked array of shape
\((16, 4, 4)\) by the helper function \texttt{\_gamma\_stack}.


% ===========================================================================
\section{Boosted Smearing}
\label{sec:boosted_smearing}
% ===========================================================================

To improve the overlap of the interpolating operators with the pion ground state
at non-zero momentum, the code employs \emph{boosted Gaussian smearing} with
momentum injection.

\subsection{Kernel}

The smearing kernel in position space is

\begin{equation}
K(\mathbf{r}; w, \mathbf{k})
  = \exp\!\Big(-\frac{\mathbf{r}^2}{2w^2}\Big) \,
    \exp\!\left[
      2\pi i
      \left(
        \frac{k_x r_x}{L_x}
        + \frac{k_y r_y}{L_y}
        + \frac{k_z r_z}{L_z}
      \right)
    \right] ,
\label{eq:kernel}
\end{equation}

where \(w\) is the Gaussian width in lattice units and
\(\mathbf{k}=(k_x,k_y,k_z)\) is the integer momentum-smearing vector used in the
code's \texttt{boost} argument.  The coordinates \(\mathbf{r}\) are centered
periodic global lattice coordinates, e.g. \(r_x\in[-L_x/2,L_x/2)\).  The first
factor is the Gaussian profile and the second factor injects the momentum
\(2\pi(k_x/L_x,k_y/L_y,k_z/L_z)\).

\subsection{Implementation}

The smearing is implemented in momentum space for efficiency.  The source
fermion is first Fourier-transformed (using \pyquda's distributed 3D FFT),
multiplied by the Fourier transform of \(K(\mathbf{r})\), and then
inverse-transformed back to position space.  This is done independently for each
spin-color component of the propagator or fermion field.

Equivalently, for a fermion field \(\psi\),
\begin{equation}
  \psi^{\,{\rm smear}}_{\mathbf{k}}(\mathbf{x},t)
  =
  \mathcal{F}^{-1}_{3d}
  \left[
    \mathcal{F}_{3d}[\psi](\mathbf{p},t)\,
    \mathcal{F}_{3d}[K_{\mathbf{k}}](\mathbf{p})
  \right](\mathbf{x},t).
\label{eq:fft_smearing}
\end{equation}
The FFT is purely spatial; the same kernel is broadcast to every Euclidean time
slice.

The smearing assumes the gauge links are in the identity gauge (or sufficiently
smoothed), so no additional gauge rotation is applied.

\subsection{Independent Boost Parameters}

Pion measurements involve both a quark and an antiquark line, which may require
different momentum injections.  The code supports four independent boost
vectors:

\begin{itemize}[nosep]
  \item \texttt{pos\_boost\_src}: momentum injected at the source for the
        forward (positive) quark propagator \(\Sq(x, x_0)\).
  \item \texttt{pos\_boost\_sink}: momentum injected at the sink for the forward
        quark propagator (relevant for the two-point function).
  \item \texttt{neg\_boost\_src}: momentum injected at the source for the
        backward (negative/antiquark) propagator.
  \item \texttt{neg\_boost\_sink}: momentum injected at the sink for the
        backward propagator; this is also applied before constructing the
        sequential source.
\end{itemize}

If the per-line source/sink boost parameters are not provided, they fall back
to the global \texttt{pos\_boost} and \texttt{neg\_boost} values.


% ===========================================================================
\section{Code Implementation}
% ===========================================================================

The contraction module resides in
\texttt{pyquda\_measurement\_utils/pion\_EMFF\_vibe\_develop.py} and defines
the class \texttt{pion\_EMFF} with two principal methods.  The sequential source
is built by \texttt{create\_meson\_bw\_seq\_pyquda} in
\texttt{pyquda\_measurement\_utils/bw\_seq\_pyquda.py}.

% ---------------------------------------------------------------------------
\subsection{The Antiquark Backward Line: \texttt{\_meson\_backward\_line}}
% ---------------------------------------------------------------------------

Given a quark propagator \(S(x, x_0)\), this function computes the antiquark
line using \(\gfive\)-hermiticity:

\begin{equation}
\Santi(x, x_0) = \gfive \, S(x, x_0)^\dagger \, \gfive .
\label{eq:mesh_backward}
\end{equation}

The corresponding \einsum{} pattern is

\begin{verbatim}
  xp.einsum("ij, wtzyxilab, kl -> wtzyxkjba",
            gamma5, prop.data.conj(), gamma5)
\end{verbatim}

In this pattern:
\begin{itemize}[nosep]
  \item \(\texttt{gamma5}_{ij}\) (index \(i\) = row, \(j\) = column) acts from
        the left on the spin\_sink index of the conjugated propagator.
  \item \(\texttt{prop.data.conj()}_{\texttt{wtzyxilab}}\) has spin\_sink \(i\)
        and spin\_src \(l\).
  \item \(\texttt{gamma5}_{kl}\) (index \(k\) = row, \(l\) = column) acts from
        the right on the spin\_src index.
  \item The output has spin\_sink \(k\), spin\_src \(j\), color\_sink \(b\),
        color\_src \(a\) --- the color indices are swapped relative to the input
        because conjugation exchanges sink and source.
\end{itemize}

% ---------------------------------------------------------------------------
\subsection{Gamma Matrix Stack: \texttt{\_gamma\_stack}}
% ---------------------------------------------------------------------------

This function constructs a stacked array of all 16 gamma matrices with shape
\((16, 4, 4)\), labeled by the index \(g\) in subsequent contractions:

\begin{equation}
\Gamma_{g;\,ij} = \texttt{gamma\_ls}[g, i, j] .
\label{eq:gamma_stack}
\end{equation}

% ---------------------------------------------------------------------------
\subsection{Source Gamma Stack: \texttt{\_source\_gamma\_stack}}
% ---------------------------------------------------------------------------

Given the sink/current gamma stack and the \texttt{src\_gamma} mode, this
function returns the corresponding source gamma stack.  For the four modes
described in Section~\ref{sec:src_gamma}:

\begin{itemize}[nosep]
  \item \texttt{"fixed\_g5"}: \(\gsrc = \gfive\) for all \(g\).
  \item \texttt{"same\_as\_sink"}: \(\gsrc^{(g)} = \gsink^{(g)}\).
  \item \texttt{"dagger\_of\_sink"}: \(\gsrc^{(g)} = \gfive
        (\gsink^{(g)})^\dagger \gfive\), implemented via\\
        \(\texttt{xp.einsum("ab, gbc, cd -> gad", gamma5,}\\
        \texttt{sink\_gamma\_ls.conj().swapaxes(1,2), gamma5)}\).
  \item Specific label: \(\gsrc^{(g)} = \Gamma_{\text{label}}\) for all \(g\).
\end{itemize}

% ---------------------------------------------------------------------------
\subsection{Two-Point Contraction: \texttt{contract\_2pt\_pion}}
% ---------------------------------------------------------------------------

The contraction proceeds in three tensor steps:

\textbf{Step 1 --- Sink gamma insertion on the backward line:}
\begin{verbatim}
  bw_prop = xp.einsum("wtzyxjicf, gim -> gwtzyxjmcf",
                       bw_prop, sink_gamma_ls)
\end{verbatim}

Here \(\texttt{bw\_prop} = \Santi(\mathbf{x}, t; \mathbf{x}_0, 0)\) with layout
\((w, t, z, y, x, \alpha, \beta, c, d)\) where \(\alpha\) = spin\_sink,
\(\beta\) = spin\_src, \(c\) = color\_sink, \(d\) = color\_src.  The sink gamma
stack has layout \((g, \alpha, \gamma)\).  The sum over \(\alpha\) inserts the
sink gamma between the backward line's spin\_src and the forward line:

\begin{equation}
\big(\Santi \cdot \gsink\big)_{\alpha\gamma; cd}
  = {\Santi}_{\alpha\beta; cd} \; (\gsink)_{\beta\gamma} .
\label{eq:sink_insert}
\end{equation}

\textbf{Step 2 --- Spin-color trace:}
\begin{verbatim}
  corr_local = xp.einsum("gwtzyxjiab, wtzyxilba, glj -> gwtzyx",
                          bw_prop, prop_pos.data, source_gamma_ls)
\end{verbatim}

The three tensors entering this contraction are:
\begin{itemize}[nosep]
  \item \(\texttt{bw\_prop}\) (after Step 1): \((\Santi \cdot \gsink)_{\alpha\beta; cd}\)
        with indices \((g, w, t, z, y, x, \alpha, \beta, c, d)\).
  \item \(\texttt{prop\_pos.data}\): \(\Sq(\mathbf{x}, t; \mathbf{x}_0, 0)\)
        with indices \((w, t, z, y, x, \beta, \gamma, d, c)\) --- note that
        the color indices are transposed relative to the backward line.
  \item \(\texttt{source\_gamma\_ls}\): \(\gsrc^{(g)}\) with indices
        \((g, \gamma, \alpha)\).
\end{itemize}

The summed indices are:
\begin{itemize}[nosep]
  \item \(\beta\) --- spin\_src of backward line with spin\_sink of forward line;
  \item \(\gamma\) --- spin\_src of forward line with spin\_sink of source gamma;
  \item \(\alpha\) --- spin\_sink of backward line with spin\_src of source gamma
        (closing the spin trace);
  \item \(c\) --- color\_sink of backward line with color\_src of forward line;
  \item \(d\) --- color\_src of backward line with color\_sink of forward line
        (closing the color trace).
\end{itemize}

The output has shape \((g, w, t, z, y, x)\) --- the correlator for each gamma
channel at each spacetime point.

In physics notation this is precisely:

\begin{equation}
\operatorname{corr\_local}_{g}(\mathbf{x}, t)
  = \tr_{\text{spin,color}}\!\Big[
      \big(\Santi \cdot \gsink^{(g)}\big)(\mathbf{x}, t) \;
      \Sq(\mathbf{x}, t) \;
      \gsrc^{(g)}
    \Big] .
\label{eq:corr_local}
\end{equation}

\textbf{Step 3 --- Momentum projection:}
\begin{verbatim}
  corr = xp.einsum("qwtzyx, gwtzyx -> gqt", phases, corr_local)
\end{verbatim}

where \(\texttt{phases}_{q, w, t, z, y, x} = e^{-i\mathbf{p}_q\cdot(\mathbf{x}-\mathbf{x}_0)}\)
for each momentum mode \(q\) in \texttt{p\_2pt}.  The sum over spatial indices
\((z, y, x)\) and the checkerboard parity \(w\) is implicit in the \einsum{}
contraction of the spacetime axes.

% ---------------------------------------------------------------------------
\subsection{Three-Point Contraction: \texttt{contract\_EMFF}}
% ---------------------------------------------------------------------------

The three-point contraction follows an identical spin-color trace structure, but
replaces the backward line from the local propagator with the sequential
antiquark propagator:

\begin{verbatim}
  seq_bw_line = _meson_backward_line(seq_bw_prop)
  current_inserted = xp.einsum(
      "wtzyxjicf, gim -> gwtzyxjmcf",
      seq_bw_line, current_gamma_ls)
  corr_local = xp.einsum(
      "gwtzyxjiab, wtzyxilba, glj -> gwtzyx",
      current_inserted, prop_pos.data, source_gamma_ls)
  corr = xp.einsum("qwtzyx, gwtzyx -> gqt", phases, corr_local)
\end{verbatim}

The correspondence with the analytic expression
Eq.~\eqref{eq:C3_seq} is:

\begin{itemize}[nosep]
  \item \(\texttt{seq\_bw\_line}\) is \(\Sseq^{\text{anti}}(x, x_0;
        \mathbf{p}_f, t_{\text{sep}})\) from Eq.~\eqref{eq:Sseq_anti}.
  \item \(\texttt{current\_gamma\_ls}\) is the stack of 16 gamma matrices for
        the current insertion \(\gcur\).
  \item The \einsum{} \(\texttt{wtzyxjicf, gim -> gwtzyxjmcf}\) inserts the
        current gamma \(\gcur\) into the sequential backward line, producing
        \(\big(\Sseq^{\text{anti}} \cdot \gcur\big)\).
  \item The second \einsum{} performs the full spin-color trace
        \(\tr[\Sseq^{\text{anti}} \cdot \gcur \cdot \Sq \cdot \gsrc]\).
  \item The final \einsum{} projects onto momentum transfer \(\mathbf{q}\) via
        Fourier phase \(e^{-i\mathbf{q}\cdot(\mathbf{x}-\mathbf{x}_0)}\).
\end{itemize}

\textbf{Result layout:} After gathering across MPI ranks, the output array has
shape \((\texttt{nq}, \texttt{ngamma}, L_t)\) or equivalently
\((g, q, t)\).  In the main application code, the axes are transposed to
\((\texttt{ngamma}, \texttt{nq}, t_{\text{sep}}+2)\) for storage.

% ---------------------------------------------------------------------------
\subsection{Sequential Source Construction}
% ---------------------------------------------------------------------------

The sequential source is built by \texttt{create\_meson\_bw\_seq\_pyquda} in
\texttt{bw\_seq\_pyquda.py}.  The steps are:

\begin{enumerate}[nosep]
  \item \textbf{Time slicing:} The input propagator \(S^{\text{neg}}(y, x_0)\)
        is restricted to time slice \(t_{\text{sink}}\) using
        \(\texttt{sequential12}(\texttt{prop}, t_{\text{sink}})\).
  \item \textbf{Momentum phase:} Multiply by
        \(e^{-i\mathbf{p}_f\cdot(\mathbf{y}-\mathbf{x}_0)}\) on the
        spatial lattice sites at the sink time slice.
        \begin{verbatim}
  seq_data = xp.einsum("wtzyx, wtzyxilab -> wtzyxilab",
                        mom_phase, seq_data)
        \end{verbatim}
  \item \textbf{Gamma insertion:} Apply \(\gseq = \gfive \gsink^\dagger \gfive\).
        \begin{verbatim}
  gamma_seq = xp.einsum("ab, cb, cd -> ad",
                         G5, sink_gamma.conj(), G5)
  seq_data = xp.einsum("ij, wtzyxjlab -> wtzyxilab",
                        gamma_seq, seq_data)
        \end{verbatim}
  \item \textbf{Boosted smearing:} The sequential source is smeared with the
        same width and boost as the original propagator,
        \(\texttt{boosted\_smearing}(\texttt{src}, w, \texttt{boost})\).
  \item \textbf{Inversion:} Solve \(D \, \Sseq = \etaseq\) via
        \(\texttt{core.invertPropagator}(\texttt{dirac}, \texttt{src\_seq}, 1, 0)\).
  \item \textbf{Final transformation:} Apply \(\gfive\) from the right to
        properly close the spin structure:
        \(\Sseq \to \Sseq \, \gfive\).  (This is a convention for the subsequent
        antiquark conversion in \texttt{contract\_EMFF}.)
        \begin{verbatim}
  final_term = xp.einsum("wtzyxijfc, ik -> wtzyxjkcf",
                          prop_smeared.data.conj(), G5)
        \end{verbatim}
\end{enumerate}

The complete sequential source is thus:

\begin{equation}
\etaseq(y; \mathbf{p}_f, t_{\text{sep}})
  = \delta_{t_y,\,t_{\text{sep}}} \;
    e^{-i\mathbf{p}_f\cdot(\mathbf{y}-\mathbf{x}_0)} \;
    \big(\gfive \, \gsink^\dagger \, \gfive\big) \;
    S^{\text{neg}}(y, x_0) .
\label{eq:etaseq_full}
\end{equation}


% ===========================================================================
\section{HDF5 Output Structure}
% ===========================================================================

The two-point and three-point correlators are saved in \hdfcmd{} format under the
data directory specified by the user (typically \texttt{./data/}).

\subsection{Two-Point Data}

\textbf{Path template:}
\begin{verbatim}
  data/c2pt/<lat>.c2pt.<cfg>.<ama>.<src_pos>.<sm_tag>.h5
\end{verbatim}

\textbf{Internal structure:}
\begin{verbatim}
  SS/
    "5"/
      PX0PY0PZ0    : array(shape=(Lt,), dtype=complex)
      PX0PY0PZ1    : ...
      ...
    "T"/
      ...
    ...
\end{verbatim}

Under the \texttt{SS} group, each gamma label contains datasets labeled by
momentum \texttt{PX<px>PY<py>PZ<pz>}.  Each dataset is a 1D complex array over
time slices, already time-rolled so that the source is at \(t=0\).  The two-point
function is saved via \(\texttt{save\_proton\_c2pt\_hdf5}\), which also parses
the source time from the tag to perform the roll.

\subsection{Three-Point (EMFF) Data}

\textbf{Path template:}
\begin{verbatim}
  data/pion_EMFF/<lat>.pion_EMFF.<cfg>.<ama>.<src_pos>.<sm_tag>.<pf_tag>.<gamma_label>.h5
\end{verbatim}

where \(\texttt{<pf\_tag>} = \texttt{PX<px>PY<py>PZ<pz>dt<t\_insert>}\)
and \(\texttt{<gamma\_label>}\) is one of the 16 gamma labels.

\textbf{Internal structure:}
\begin{verbatim}
  SS/
    "5"/
      PX0PY0PZ0    : array(shape=(tsep+2,), dtype=complex)
      PX1PY0PZ0    : ...
      ...
\end{verbatim}

Under the \texttt{SS} group, the single gamma label (the file contains one gamma
channel per file after MPI-distributed writing) contains datasets labeled by
momentum transfer \texttt{PX<qx>PY<qy>PZ<qz>}.  Each dataset is a 1D complex array
of length \(t_{\text{sep}}+2\), indexed by the current insertion time slice
\(\tau\).

The saving is performed by \(\texttt{save\_pion\_EMFF\_hdf5\_noRoll}\), which
does not apply any time roll (the source-time alignment is handled during
post-processing).

\subsection{File Naming Convention Summary}

\begin{center}
\begin{tabular}{ll}
\toprule
\textbf{Component} & \textbf{Description} \\
\midrule
\(\texttt{<lat>}\)        & Lattice tag (e.g., \texttt{S8T32}) \\
\(\texttt{<cfg>}\)        & Configuration number \\
\(\texttt{<ama>}\)        & AMA label (e.g., \texttt{EMFF.ex}) \\
\(\texttt{<src\_pos>}\)   & Source position
                          \texttt{x<x>y<y>z<z>t<t>} \\
\(\texttt{<sm\_tag>}\)     & Smearing tag (HYP steps, source type, width,
                          boost vectors) \\
\(\texttt{<pf\_tag>}\)     & Final momentum and time separation \\
\(\texttt{<gamma\_label>}\) & Gamma matrix label (5, T, T5, ...) \\
\bottomrule
\end{tabular}
\end{center}

The tag-building helpers are defined in
\texttt{pyquda\_measurement\_utils/io\_corr.py}:
\begin{itemize}[nosep]
  \item \(\texttt{get\_c2pt\_file\_tag}\) --- builds the two-point HDF5 path.
  \item \(\texttt{get\_pion\_EMFF\_file\_tag}\) --- builds the three-point
        HDF5 path.
  \item \(\texttt{save\_proton\_c2pt\_hdf5}\) --- writes two-point data
        (shared with baryon workflows).
  \item \(\texttt{save\_pion\_EMFF\_hdf5\_noRoll}\) --- writes three-point
        EMFF data.
\end{itemize}


% ===========================================================================
\section{Workflow Summary}
% ===========================================================================

The complete pion EMFF measurement, as orchestrated by the main application
script \texttt{Pyquda\_pion\_EMFF.py}, proceeds as follows:

\begin{enumerate}[nosep]
  \item \textbf{Gauge field preparation:} Read the gauge configuration, apply
        HYP smearing (1 step, \(\alpha = (0.75, 0.6, 0.3)\)), set antiperiodic
        temporal boundary conditions.
  \item \textbf{Dirac operator:} Construct a Wilson-clover Dirac operator with
        the specified mass, \(c_{\text{SW}}\), and multigrid parameters.
  \item \textbf{Source creation:} Generate point sources at distributed
        spacetime positions.
  \item \textbf{Positive quark propagator:} Apply boosted smearing at the source
        (\(\texttt{pos\_boost\_src}\)), invert.
  \item \textbf{Negative quark propagator:} Apply boosted smearing at the source
        (\(\texttt{neg\_boost\_src}\)) if different from positive; otherwise
        copy the positive propagator.
  \item \textbf{Two-point contraction:} Apply sink smearing to both propagators,
        contract via \(\texttt{contract\_2pt\_pion}\), save to
        \(\texttt{data/c2pt/}\).
  \item \textbf{Sequential propagator:} Apply sink smearing to the negative
        propagator (\(\texttt{neg\_boost\_sink}\)), build the sequential source
        at \(t_{\text{sink}}\), invert, and post-process.
  \item \textbf{Three-point contraction:} Contract the sequential propagator
        with the positive forward propagator and all 16 current gamma
        structures, project onto momentum transfer, save per-gamma-channel
        files to \(\texttt{data/pion\_EMFF/}\).
\end{enumerate}


% ===========================================================================
\section{Code Source Reference}
% ===========================================================================

The primary source files comprising the pion EMFF measurement are:

\begin{itemize}[nosep]
  \item \texttt{pyquda\_measurement\_utils/pion\_EMFF\_vibe\_develop.py}
        --- Defines the \texttt{pion\_EMFF} class with
        \(\texttt{contract\_2pt\_pion}\) and \(\texttt{contract\_EMFF}\),
        the gamma matrix stacks, the antiquark backward line conversion, and
        the source gamma mode logic.

  \item \texttt{pyquda\_measurement\_utils/bw\_seq\_pyquda.py}
        --- Defines \(\texttt{create\_meson\_bw\_seq\_pyquda}\) for building
        the fixed-sink meson sequential source, time slicing, and inversion.
        Also contains \(\texttt{create\_bw\_seq\_pyquda}\) for baryon
        sequential sources (used in other workflows).

  \item \texttt{pyquda\_measurement\_utils/boosted\_smearing\_pyquda.py}
        --- Implements the FFT boosted Gaussian smearing kernel with momentum
        injection via distributed 3D FFT.

  \item \texttt{pyquda\_measurement\_utils/io\_corr.py}
        --- HDF5 I/O helpers for two-point, three-point, and other
        correlator output formats.

  \item \texttt{pyquda\_measurement\_utils/tools.py}
        --- Backend-agnostic utilities (\(\texttt{\_get\_xp\_from\_array}\),
        \(\texttt{\_asarray\_on\_queue}\), \(\texttt{mpi\_print}\)).

  \item \texttt{application/EMFF\_pion/perlmutter/Pyquda\_pion\_EMFF.py}
        --- Main application script that orchestrates the full measurement
        workflow, handling gauge I/O, propagator generation, and HDF5 output.
\end{itemize}

The sequential source time-slicing utility \(\texttt{sequential12}\) is provided
by \(\texttt{pyquda\_utils.source}\).  The momentum phase factors are generated
by \(\texttt{pyquda\_utils.phase.MomentumPhase}\).

\end{document}
